\documentclass[aspectratio=169]{beamer}
%--- Configuración del Tema y Colores ---
\usetheme{Madrid}
\usecolortheme{whale}
\setbeamertemplate{navigation symbols}{}
\setbeamertemplate{caption}[numbered]

%--- Paquetes ---
\usepackage[utf8]{inputenc}
\usepackage[spanish]{babel}
\usepackage{graphicx}
\usepackage{booktabs}
\usepackage{tikz}
\usepackage{listings}
\usepackage{xcolor}

%--- Configuración de Listings para Python ---
\definecolor{codegreen}{rgb}{0,0.6,0}
\definecolor{codegray}{rgb}{0.5,0.5,0.5}
\definecolor{codepurple}{rgb}{0.58,0,0.82}
\definecolor{backcolour}{rgb}{0.95,0.95,0.92}

\lstdefinestyle{pythonstyle}{
    backgroundcolor=\color{backcolour},
    commentstyle=\color{codegreen},
    keywordstyle=\color{blue}\bfseries,
    numberstyle=\tiny\color{codegray},
    stringstyle=\color{codepurple},
    basicstyle=\ttfamily\scriptsize,
    breakatwhitespace=false,
    breaklines=true,
    captionpos=b,
    keepspaces=true,
    numbers=left,
    numbersep=5pt,
    showspaces=false,
    showstringspaces=false,
    showtabs=false,
    tabsize=2,
    language=Python,
    morekeywords={self, True, False, None, class, def, return, if, else, elif, for, while, in, and, or, not, lambda, try, except, finally, raise, with, as, from, import, pass, break, continue, global, nonlocal, assert, yield, del, is, async, await}
}
\lstset{style=pythonstyle}

%--- Metadatos del Documento ---
\title[PM2]{Programación Matemática 2}
\subtitle{Clase 3: Paradigma Orientado a Objetos}
\author[Luis Alfredo Alvarado]{Prof. Luis Alfredo Alvarado Rodríguez}
\institute[ECFM - USAC]{
    Escuela de Ciencias Físicas y Matemáticas \\
    Universidad de San Carlos de Guatemala
}
\date{Primer Semestre, 2026}

\begin{document}

%--- Diapositiva de Título ---
\begin{frame}
    \titlepage
\end{frame}

%--- Diapositiva: Tabla de Contenidos ---
\begin{frame}{Agenda de Hoy}
    \tableofcontents
\end{frame}

%===========================================================
% SECCIÓN 1: INTRODUCCIÓN
%===========================================================
\section{Introducción a la Programación Orientada a Objetos}

\begin{frame}{¿Qué es la Programación Orientada a Objetos?}
    \begin{block}{Definición}
        La \textbf{Programación Orientada a Objetos (POO)} es un paradigma que organiza el software como una colección de \textbf{objetos} que contienen datos y comportamiento.
    \end{block}
    \vspace{0.5cm}
    
    \begin{alertblock}{Filosofía}
        \textbf{``Se modela el mundo real mediante entidades que tienen características (atributos) y pueden realizar acciones (métodos).''}
    \end{alertblock}
    \vspace{0.3cm}
    
    \textbf{Analogía:} Un automóvil tiene características (color, marca, modelo) y puede realizar acciones (acelerar, frenar, girar).
\end{frame}

\begin{frame}{Los Cuatro Pilares de la POO}
    \begin{columns}[t]
        \column{0.48\textwidth}
        \begin{block}{1. Encapsulamiento}
            Se ocultan los detalles internos y se expone solo lo necesario.
        \end{block}
        \vspace{0.3cm}
        \begin{block}{2. Abstracción}
            Se simplifica la complejidad mostrando solo las características esenciales.
        \end{block}
        
        \column{0.48\textwidth}
        \begin{block}{3. Herencia}
            Se crean nuevas clases basadas en clases existentes.
        \end{block}
        \vspace{0.3cm}
        \begin{block}{4. Polimorfismo}
            Se permite que objetos de diferentes clases respondan al mismo mensaje.
        \end{block}
    \end{columns}
\end{frame}

%===========================================================
% SECCIÓN 2: CLASES Y OBJETOS
%===========================================================
\section{Clases y Objetos}

\begin{frame}{Clases y Objetos: Conceptos Fundamentales}
    \begin{columns}
        \column{0.5\textwidth}
        \begin{exampleblock}{Clase}
            Es el \textbf{molde} o \textbf{plantilla} que define:
            \begin{itemize}
                \item Atributos (datos)
                \item Métodos (comportamiento)
            \end{itemize}
        \end{exampleblock}
        
        \column{0.5\textwidth}
        \begin{alertblock}{Objeto}
            Es una \textbf{instancia} de una clase:
            \begin{itemize}
                \item Entidad concreta en memoria
                \item Posee valores específicos
            \end{itemize}
        \end{alertblock}
    \end{columns}
    
    \vspace{0.5cm}
    \centering
    \begin{tikzpicture}
        \node[draw, rectangle, fill=blue!20, minimum width=2.5cm, minimum height=1cm] (clase) {Clase: Perro};
        \node[draw, rectangle, fill=green!20, minimum width=2cm, minimum height=0.8cm, right=1cm of clase] (obj1) {firulais};
        \node[draw, rectangle, fill=green!20, minimum width=2cm, minimum height=0.8cm, right=0.5cm of obj1] (obj2) {max};
        \node[draw, rectangle, fill=green!20, minimum width=2cm, minimum height=0.8cm, right=0.5cm of obj2] (obj3) {luna};
        \draw[->, thick] (clase) -- (obj1);
        \draw[->, thick] (clase) -- (obj2);
        \draw[->, thick] (clase) -- (obj3);
        \node[below=0.1cm of clase] {\scriptsize Plantilla};
        \node[below=0.1cm of obj2] {\scriptsize Instancias (Objetos)};
    \end{tikzpicture}
\end{frame}

\begin{frame}[fragile]{Definición de una Clase en Python}
    \begin{lstlisting}[title=Sintaxis Básica de una Clase]
class Persona:
    """Clase que representa a una persona."""
    
    # Constructor: se ejecuta al crear un objeto
    def __init__(self, nombre, edad):
        # Atributos de instancia
        self.nombre = nombre
        self.edad = edad
    
    # Metodo de instancia
    def saludar(self):
        return f"Hola, mi nombre es {self.nombre}"
    
    # Metodo de instancia
    def es_mayor_de_edad(self):
        return self.edad >= 18
    \end{lstlisting}
    
    \textbf{Nótese:} \texttt{self} hace referencia al objeto actual (instancia).
\end{frame}

\begin{frame}[fragile]{Creación y Uso de Objetos}
    \begin{lstlisting}[title=Instanciación de Objetos]
# Se crean objetos (instancias) de la clase Persona
persona1 = Persona("Ana", 25)
persona2 = Persona("Carlos", 17)

# Se accede a los atributos
print(persona1.nombre)  # Ana
print(persona2.edad)    # 17

# Se invocan los metodos
print(persona1.saludar())           # Hola, mi nombre es Ana
print(persona1.es_mayor_de_edad())  # True
print(persona2.es_mayor_de_edad())  # False

# Cada objeto tiene su propio estado
persona1.edad = 26  # Se modifica solo persona1
print(persona1.edad)  # 26
print(persona2.edad)  # 17 (sin cambios)
    \end{lstlisting}
\end{frame}

\begin{frame}[fragile]{Atributos de Clase vs Atributos de Instancia}
    \begin{lstlisting}[title=Tipos de Atributos]
class Estudiante:
    # Atributo de CLASE (compartido por todas las instancias)
    universidad = "USAC"
    contador = 0
    
    def __init__(self, nombre, carnet):
        # Atributos de INSTANCIA (unicos por objeto)
        self.nombre = nombre
        self.carnet = carnet
        Estudiante.contador += 1  # Se incrementa el contador

# Se crean instancias
est1 = Estudiante("Maria", "202012345")
est2 = Estudiante("Pedro", "202054321")

print(est1.universidad)      # USAC (atributo de clase)
print(est2.universidad)      # USAC (mismo valor)
print(Estudiante.contador)   # 2 (compartido)
    \end{lstlisting}
\end{frame}

%===========================================================
% SECCIÓN 3: ENCAPSULAMIENTO
%===========================================================
\section{Encapsulamiento}

\begin{frame}{Encapsulamiento: Concepto}
    \begin{block}{Definición}
        El \textbf{encapsulamiento} consiste en ocultar los detalles internos de un objeto y exponer solo una interfaz pública controlada.
    \end{block}
    \vspace{0.5cm}
    
    \textbf{Convenciones en Python:}
    \begin{itemize}
        \item \texttt{atributo} — Público (accesible desde cualquier lugar)
        \item \texttt{\_atributo} — Protegido (convención: uso interno)
        \item \texttt{\_\_atributo} — Privado (name mangling)
    \end{itemize}
    \vspace{0.3cm}
    
    \textbf{Beneficios:}
    \begin{itemize}
        \item Se protege la integridad de los datos.
        \item Se reduce el acoplamiento entre componentes.
        \item Se facilita el mantenimiento del código.
    \end{itemize}
\end{frame}

\begin{frame}[fragile]{Ejemplo de Encapsulamiento}
    \begin{lstlisting}[title=Atributos Privados y Métodos de Acceso]
class CuentaBancaria:
    def __init__(self, titular, saldo_inicial):
        self.titular = titular
        self.__saldo = saldo_inicial  # Atributo privado
    
    # Getter: permite leer el saldo
    def obtener_saldo(self):
        return self.__saldo
    
    # Metodos que controlan el acceso al saldo
    def depositar(self, monto):
        if monto > 0:
            self.__saldo += monto
            return True
        return False
    
    def retirar(self, monto):
        if 0 < monto <= self.__saldo:
            self.__saldo -= monto
            return True
        return False
    \end{lstlisting}
\end{frame}

\begin{frame}[fragile]{Uso del Encapsulamiento}
    \begin{lstlisting}[title=Protección de Datos]
cuenta = CuentaBancaria("Juan Perez", 1000)

# Acceso controlado mediante metodos
print(cuenta.obtener_saldo())  # 1000
cuenta.depositar(500)
print(cuenta.obtener_saldo())  # 1500

cuenta.retirar(200)
print(cuenta.obtener_saldo())  # 1300

# Intento de acceso directo (falla)
# print(cuenta.__saldo)  # AttributeError

# El saldo no puede ser negativo gracias al encapsulamiento
cuenta.retirar(5000)  # Retorna False, no se ejecuta
print(cuenta.obtener_saldo())  # 1300 (sin cambios)
    \end{lstlisting}
    
    \textbf{Nótese:} La lógica de negocio se protege dentro de la clase.
\end{frame}

\begin{frame}[fragile]{Properties: Encapsulamiento Pythónico}
    \begin{lstlisting}[title=Uso de @property]
class Rectangulo:
    def __init__(self, ancho, alto):
        self._ancho = ancho
        self._alto = alto
    
    @property
    def area(self):
        """Se calcula automaticamente (solo lectura)."""
        return self._ancho * self._alto
    
    @property
    def ancho(self):
        return self._ancho
    
    @ancho.setter
    def ancho(self, valor):
        if valor > 0:
            self._ancho = valor
        else:
            raise ValueError("El ancho debe ser positivo")
    \end{lstlisting}
\end{frame}

\begin{frame}[fragile]{Uso de Properties}
    \begin{lstlisting}[title=Acceso Transparente]
rect = Rectangulo(10, 5)

# Se accede como atributo, pero ejecuta el metodo
print(rect.area)   # 50
print(rect.ancho)  # 10

# Se usa el setter automaticamente
rect.ancho = 20
print(rect.area)   # 100

# Validacion integrada
# rect.ancho = -5  # ValueError: El ancho debe ser positivo
    \end{lstlisting}
    
    \vspace{0.3cm}
    \textbf{Ventaja:} Se mantiene la sintaxis simple de atributos con la protección de métodos.
\end{frame}

%===========================================================
% SECCIÓN 4: HERENCIA
%===========================================================
\section{Herencia}

\begin{frame}{Herencia: Concepto}
    \begin{block}{Definición}
        La \textbf{herencia} permite crear nuevas clases (hijas) que reutilizan, extienden o modifican el comportamiento de clases existentes (padres).
    \end{block}
    \vspace{0.3cm}
    
    \centering
    \begin{tikzpicture}[node distance=1.5cm]
        \node[draw, rectangle, fill=blue!20, minimum width=3cm, minimum height=1cm] (animal) {Animal};
        \node[draw, rectangle, fill=green!20, minimum width=2.5cm, minimum height=0.8cm, below left=1cm and 0.5cm of animal] (perro) {Perro};
        \node[draw, rectangle, fill=green!20, minimum width=2.5cm, minimum height=0.8cm, below right=1cm and 0.5cm of animal] (gato) {Gato};
        \draw[->, thick] (perro) -- (animal);
        \draw[->, thick] (gato) -- (animal);
        \node[right=0.3cm of animal] {\scriptsize Clase Padre (Base)};
        \node[left=0.3cm of perro] {\scriptsize Clases Hijas};
    \end{tikzpicture}
    
    \vspace{0.3cm}
    \textbf{Relación:} ``Es un'' (un Perro \textit{es un} Animal)
\end{frame}

\begin{frame}[fragile]{Ejemplo de Herencia}
    \begin{lstlisting}[title=Clase Base]
class Animal:
    def __init__(self, nombre, edad):
        self.nombre = nombre
        self.edad = edad
    
    def presentarse(self):
        return f"Soy {self.nombre} y tengo {self.edad} anios"
    
    def hacer_sonido(self):
        return "..."  # Sonido generico
    \end{lstlisting}
\end{frame}

\begin{frame}[fragile]{Clases Derivadas}
    \begin{lstlisting}[title=Clases Hijas que Heredan de Animal]
class Perro(Animal):  # Hereda de Animal
    def __init__(self, nombre, edad, raza):
        super().__init__(nombre, edad)  # Llama al constructor padre
        self.raza = raza
    
    def hacer_sonido(self):  # Sobreescribe el metodo
        return "Guau guau!"
    
    def traer_pelota(self):  # Metodo exclusivo de Perro
        return f"{self.nombre} trae la pelota"

class Gato(Animal):
    def __init__(self, nombre, edad, color):
        super().__init__(nombre, edad)
        self.color = color
    
    def hacer_sonido(self):
        return "Miau!"
    \end{lstlisting}
\end{frame}

\begin{frame}[fragile]{Uso de la Herencia}
    \begin{lstlisting}[title=Instanciación y Uso]
# Se crean instancias de las clases hijas
perro = Perro("Firulais", 3, "Labrador")
gato = Gato("Michi", 2, "Naranja")

# Metodos heredados
print(perro.presentarse())  # Soy Firulais y tengo 3 anios
print(gato.presentarse())   # Soy Michi y tengo 2 anios

# Metodos sobreescritos
print(perro.hacer_sonido())  # Guau guau!
print(gato.hacer_sonido())   # Miau!

# Metodo exclusivo
print(perro.traer_pelota())  # Firulais trae la pelota

# Verificacion de herencia
print(isinstance(perro, Animal))  # True
print(isinstance(perro, Perro))   # True
    \end{lstlisting}
\end{frame}

\begin{frame}[fragile]{Herencia Múltiple}
    \begin{lstlisting}[title=Una Clase Puede Heredar de Múltiples Padres]
class Volador:
    def volar(self):
        return "Estoy volando"

class Nadador:
    def nadar(self):
        return "Estoy nadando"

class Pato(Animal, Volador, Nadador):
    def hacer_sonido(self):
        return "Cuac cuac!"

# El pato hereda de todas las clases
pato = Pato("Donald", 5)
print(pato.presentarse())   # Heredado de Animal
print(pato.volar())         # Heredado de Volador
print(pato.nadar())         # Heredado de Nadador
print(pato.hacer_sonido())  # Sobreescrito en Pato
    \end{lstlisting}
\end{frame}

%===========================================================
% SECCIÓN 5: POLIMORFISMO
%===========================================================
\section{Polimorfismo}

\begin{frame}{Polimorfismo: Concepto}
    \begin{block}{Definición}
        El \textbf{polimorfismo} permite que objetos de diferentes clases respondan al mismo mensaje (método) de formas distintas.
    \end{block}
    \vspace{0.3cm}
    
    \textbf{Del griego:} \textit{poly} (muchas) + \textit{morphe} (formas)
    \vspace{0.5cm}
    
    \textbf{Tipos en Python:}
    \begin{itemize}
        \item \textbf{Polimorfismo de subtipo:} Clases hijas redefinen métodos de la clase padre.
        \item \textbf{Duck Typing:} ``Si camina como pato y hace cuac como pato, entonces es un pato.''
    \end{itemize}
    \vspace{0.3cm}
    
    \textbf{Beneficio:} Se puede escribir código genérico que funciona con múltiples tipos.
\end{frame}

\begin{frame}[fragile]{Ejemplo de Polimorfismo}
    \begin{lstlisting}[title=Polimorfismo con Herencia]
def hacer_hablar(animal):
    """Funcion que funciona con cualquier Animal."""
    print(f"{animal.nombre} dice: {animal.hacer_sonido()}")

# Se crean diferentes tipos de animales
animales = [
    Perro("Rex", 4, "Pastor Aleman"),
    Gato("Luna", 2, "Blanco"),
    Pato("Lucas", 1)
]

# La misma funcion funciona con todos
for animal in animales:
    hacer_hablar(animal)

# Salida:
# Rex dice: Guau guau!
# Luna dice: Miau!
# Lucas dice: Cuac cuac!
    \end{lstlisting}
\end{frame}

\begin{frame}[fragile]{Duck Typing en Python}
    \begin{lstlisting}[title=No Importa el Tipo, Importa el Comportamiento]
class Coche:
    def __init__(self, marca):
        self.marca = marca
    
    def mover(self):
        return f"El coche {self.marca} avanza por la carretera"

class Avion:
    def __init__(self, modelo):
        self.modelo = modelo
    
    def mover(self):
        return f"El avion {self.modelo} vuela por el cielo"

# Funcion generica que no verifica tipos
def iniciar_viaje(vehiculo):
    print(vehiculo.mover())

iniciar_viaje(Coche("Toyota"))  # El coche Toyota avanza...
iniciar_viaje(Avion("Boeing"))  # El avion Boeing vuela...
    \end{lstlisting}
\end{frame}

\begin{frame}[fragile]{Métodos Especiales (Dunder Methods)}
    \begin{lstlisting}[title=Polimorfismo con Operadores]
class Vector:
    def __init__(self, x, y):
        self.x = x
        self.y = y
    
    def __add__(self, otro):  # Define el operador +
        return Vector(self.x + otro.x, self.y + otro.y)
    
    def __str__(self):  # Define como se imprime
        return f"Vector({self.x}, {self.y})"
    
    def __eq__(self, otro):  # Define el operador ==
        return self.x == otro.x and self.y == otro.y

v1 = Vector(1, 2)
v2 = Vector(3, 4)
v3 = v1 + v2  # Usa __add__
print(v3)     # Vector(4, 6) - Usa __str__
    \end{lstlisting}
\end{frame}

%===========================================================
% SECCIÓN 6: ABSTRACCIÓN
%===========================================================
\section{Abstracción}

\begin{frame}{Abstracción: Concepto}
    \begin{block}{Definición}
        La \textbf{abstracción} consiste en definir una interfaz común sin implementar los detalles, obligando a las clases hijas a proporcionar la implementación.
    \end{block}
    \vspace{0.5cm}
    
    \textbf{Clases Abstractas:}
    \begin{itemize}
        \item No se pueden instanciar directamente.
        \item Definen métodos que \textbf{deben} ser implementados por las subclases.
        \item Se utiliza el módulo \texttt{abc} (Abstract Base Classes).
    \end{itemize}
    \vspace{0.3cm}
    
    \textbf{Propósito:} Establecer un ``contrato'' que las clases hijas deben cumplir.
\end{frame}

\begin{frame}[fragile]{Ejemplo de Clase Abstracta}
    \begin{lstlisting}[title=Definición de Clase Abstracta]
from abc import ABC, abstractmethod

class Figura(ABC):  # Hereda de ABC (Abstract Base Class)
    
    @abstractmethod
    def area(self):
        """Todas las figuras deben calcular su area."""
        pass
    
    @abstractmethod
    def perimetro(self):
        """Todas las figuras deben calcular su perimetro."""
        pass
    
    def descripcion(self):
        """Metodo concreto (con implementacion)."""
        return f"Area: {self.area()}, Perimetro: {self.perimetro()}"
    \end{lstlisting}
    
    \textbf{Nótese:} No se puede crear una instancia de \texttt{Figura} directamente.
\end{frame}

\begin{frame}[fragile]{Implementación de Clases Concretas}
    \begin{lstlisting}[title=Clases que Implementan la Abstracción]
class Circulo(Figura):
    def __init__(self, radio):
        self.radio = radio
    
    def area(self):
        return 3.14159 * self.radio ** 2
    
    def perimetro(self):
        return 2 * 3.14159 * self.radio

class Rectangulo(Figura):
    def __init__(self, ancho, alto):
        self.ancho = ancho
        self.alto = alto
    
    def area(self):
        return self.ancho * self.alto
    
    def perimetro(self):
        return 2 * (self.ancho + self.alto)
    \end{lstlisting}
\end{frame}

\begin{frame}[fragile]{Uso de la Abstracción}
    \begin{lstlisting}[title=Polimorfismo con Clases Abstractas]
# figura = Figura()  # TypeError: Can't instantiate abstract class

figuras = [
    Circulo(5),
    Rectangulo(4, 6),
    Circulo(3)
]

for figura in figuras:
    print(figura.descripcion())

# Salida:
# Area: 78.53975, Perimetro: 31.4159
# Area: 24, Perimetro: 20
# Area: 28.27431, Perimetro: 18.84954
    \end{lstlisting}
    
    \vspace{0.3cm}
    \textbf{Ventaja:} Se garantiza que todas las figuras tengan \texttt{area()} y \texttt{perimetro()}.
\end{frame}

%===========================================================
% SECCIÓN 7: COMPARACIÓN DE PARADIGMAS
%===========================================================
\section{Comparación de Paradigmas}

\begin{frame}{Imperativo vs Funcional vs POO}
    \begin{table}[]
        \centering
        \renewcommand{\arraystretch}{1.3}
        \scriptsize
        \begin{tabular}{p{2.5cm} p{3.5cm} p{3.5cm} p{4cm}}
            \toprule
            \textbf{Aspecto} & \textbf{Imperativo} & \textbf{Funcional} & \textbf{POO} \\
            \midrule
            Unidad básica & Instrucción & Función & Objeto \\
            Estado & Mutable & Inmutable & Encapsulado en objetos \\
            Enfoque & Cómo hacerlo & Qué resultado & Modelar entidades \\
            Reutilización & Funciones/procedimientos & Composición & Herencia/composición \\
            Ideal para & Scripts, algoritmos & Transformación de datos & Sistemas complejos \\
            \bottomrule
        \end{tabular}
    \end{table}
    
    \vspace{0.5cm}
    \centering
    \textbf{Python permite combinar los tres paradigmas según las necesidades del problema.}
\end{frame}

\begin{frame}[fragile]{Un Problema, Tres Enfoques}
    \textbf{Problema:} Calcular el área total de una lista de rectángulos.
    
    \begin{columns}
        \column{0.32\textwidth}
        \begin{lstlisting}[title=Imperativo,basicstyle=\ttfamily\tiny]
rectangulos = [
  (4, 5), (3, 7), (2, 8)
]
total = 0
for r in rectangulos:
    area = r[0] * r[1]
    total += area
print(total)  # 55
        \end{lstlisting}
        
        \column{0.32\textwidth}
        \begin{lstlisting}[title=Funcional,basicstyle=\ttfamily\tiny]
rectangulos = [
  (4, 5), (3, 7), (2, 8)
]
total = sum(
    map(
      lambda r: r[0]*r[1],
      rectangulos
    )
)
print(total)  # 55
        \end{lstlisting}
        
        \column{0.36\textwidth}
        \begin{lstlisting}[title=POO,basicstyle=\ttfamily\tiny]
class Rect:
  def __init__(self, a, h):
    self.ancho = a
    self.alto = h
  def area(self):
    return self.ancho*self.alto

rects = [Rect(4,5),
         Rect(3,7),
         Rect(2,8)]
total = sum(r.area() 
            for r in rects)
print(total)  # 55
        \end{lstlisting}
    \end{columns}
\end{frame}

%===========================================================
% SECCIÓN 8: EJERCICIOS
%===========================================================
\section{Ejercicios Propuestos}

\begin{frame}{Ejercicios para Practicar}
    \begin{enumerate}
        \setlength\itemsep{0.7em}
        \item Crear una clase \texttt{Libro} con atributos (título, autor, páginas) y un método que indique si es un libro largo ($>$ 300 páginas).
        
        \item Implementar una clase \texttt{CuentaAhorro} que herede de \texttt{CuentaBancaria} y agregue un método para calcular intereses.
        
        \item Crear una jerarquía de clases para un sistema de empleados: \texttt{Empleado} (base), \texttt{Gerente} y \texttt{Desarrollador} (hijas), cada una con su método \texttt{calcular\_salario()}.
        
        \item Implementar las clases \texttt{Triangulo} y \texttt{Cuadrado} que hereden de la clase abstracta \texttt{Figura}.
        
        \item \textbf{(Bonus)} Crear una clase \texttt{Fraccion} con métodos especiales para sumar, restar, multiplicar y comparar fracciones usando los operadores \texttt{+}, \texttt{-}, \texttt{*}, \texttt{==}.
    \end{enumerate}
\end{frame}

%--- Diapositiva Final ---
\begin{frame}
    \centering
    \Huge \textbf{¿Preguntas?}
    
    \vspace{0.8cm}
    \large
    \textbf{Próxima clase:} Introducción al Vibe Coding
    
    \vspace{0.5cm}
    \normalsize
    Prof. Luis Alfredo Alvarado Rodríguez\\
    \small Escuela de Ciencias Físicas y Matemáticas
    
    \vspace{0.5cm}
    \footnotesize
    \textit{``La programación orientada a objetos ofrece una forma sostenible de escribir código espagueti.''} — Paul Graham
\end{frame}

\end{document}